% Generated by scripts/build_baseline_report.py; compile with XeLaTeX twice.
\documentclass[UTF8,11pt,a4paper,fontset=none]{ctexart}
\usepackage[margin=18mm,top=20mm,bottom=20mm]{geometry}
\usepackage{fontspec,booktabs,longtable,array,enumitem,fancyhdr,amsmath}
\usepackage[hidelinks,unicode]{hyperref}
\IfFontExistsTF{Microsoft YaHei}{\setCJKmainfont{Microsoft YaHei}\setCJKsansfont{Microsoft YaHei}}{\setCJKmainfont{Noto Sans CJK SC}\setCJKsansfont{Noto Sans CJK SC}}
\IfFontExistsTF{Arial}{\setmainfont{Arial}}{\setmainfont{TeX Gyre Heros}}
\setmonofont{Latin Modern Mono}
\setlength{\parindent}{0pt}\setlength{\parskip}{0.65em}
\renewcommand{\arraystretch}{1.18}
\newcolumntype{L}[1]{>{\raggedright\arraybackslash}p{#1}}
\ctexset{section={format=\Large\bfseries,beforeskip=0pt,afterskip=0.7em},subsection={format=\normalsize\bfseries,beforeskip=0.7em,afterskip=0.25em}}
\pagestyle{fancy}\fancyhf{}\fancyfoot[C]{\thepage}\renewcommand{\headrulewidth}{0pt}
\setlength{\emergencystretch}{2em}
\hypersetup{pdftitle={Baseline Report},pdfauthor={},pdfsubject={复现总表、2026评测核查与非饱和数据集}}
\begin{document}
\begin{titlepage}\thispagestyle{empty}\vspace*{0.34\textheight}
\begin{center}{\fontsize{30}{42}\selectfont\bfseries Baseline Report}\end{center}
\end{titlepage}
\setcounter{page}{1}

\clearpage\section{结论与阅读规则}
本报告以一张五列总表汇总既有实验，并据 2026 年相关顶会论文确定后续评测覆盖范围。实验记录截至 2026-09-05，文献及报告核查截至 2026-09-13；本次不新增训练、策略评测或数据采集。

核心判断：标准 LIBERO 作为正确性与性能保持对照，不承担主要提升证据；LIBERO-Plus 承担扰动下的主动验证与恢复；MIKASA-Robo 的记忆任务承担独立仿真的时序证据评测；SAFE-DROID 承担真实轨迹上的外部失败检测。MIKASA-Robo 是本次选出的待接入评测，不是已经复现完成的数据集。

“2026 顶会大多需要三个数据集”没有得到本次样本支持。8 篇相关主会论文在统一口径下使用 2-3 个仿真 benchmark：5 篇为 2 个、3 篇为 3 个，并各有真机实验。该小规模定向核查不是全年论文普查，更不是投稿硬性要求。

harness 在本报告中指模型外部的执行与评测机制：输入封装、历史管理、工具调用、阈值、动作循环、恢复及成本记账。模型本身的记忆或注意力不等于外部主动证据获取。当前实验中没有已完成的“预算路由＋代码迭代”新方法成绩。

SR 与 accuracy 用百分数；AUROC / AUPRC 用 0-100 尺度，不是任务成功率。本次 SAFE 的 \(\pm\) 为 3 种子均值与样本标准差，论文 \(\pm\) 沿用原文；WCM 保留原始尺度。SAFE 主表取 early-stop 窗口内最大风险分数，成败轨迹均截到每任务最短轨迹长度；该指标不等于实际检测提前量或恢复收益。

NR 表示未在核查论文中定位到对应数据域/协议的报告值；N/A 表示该行根本不是论文模型性能指标。缺项不填 0，不从柱状图目测编造精确值。

左右两列并排用于追溯，不自动表示公平的同协议复现：OFT/AVA 的试验规模不同；SAFE 的超参数搜索范围不同；Plus 的旧诊断协议存在提示词混杂。所有这些差异在 harness 与复现列中显式披露。

\clearpage\section{Baseline 五列总表}
\begingroup\fontsize{9.2}{12.2}\selectfont\setlength{\tabcolsep}{2.5pt}
\begin{longtable}{@{}L{22mm}L{25mm}L{39mm}L{38mm}L{41mm}@{}}
\caption{既有实验总表：论文报告与本项目实际运行分别列示。}\\\toprule
{\bfseries dataset} & {\bfseries model} & {\bfseries harness} & {\bfseries reported performance\par (in paper)} & {\bfseries our reproduced performance} \\\midrule
\endfirsthead
\multicolumn{5}{l}{\textbf{表 1（续）}}\\\toprule
{\bfseries dataset} & {\bfseries model} & {\bfseries harness} & {\bfseries reported performance\par (in paper)} & {\bfseries our reproduced performance} \\\midrule
\endhead
\midrule\multicolumn{5}{r}{续下页}\\\endfoot
\bottomrule\endlastfoot
\multicolumn{5}{@{}l}{\textbf{A. 策略闭环：有限规模}}\\*[3pt]
LIBERO\par Spatial & OpenVLA-OFT & H1：作者策略循环；双相机＋本体；10 任务\(\times\)5 初态；无外加验证/恢复。 & SR 97.6\%；Spatial 专用策略。\hyperlink{src:OFT}{[1]} & SR 100.00\%（50/50）；每任务 5 次，不是论文 50 次。 \\[6pt]
LIBERO\par Spatial & AVA-VLA & H1：作者策略循环；双相机＋本体；10 任务\(\times\)5 初态；无外加验证/恢复。 & SR 97.4\%；四套件联合策略。\hyperlink{src:AVA}{[2]} & SR 96.00\%（48/50）；每任务 5 次，不是论文 50 次。 \\[6pt]
\multicolumn{5}{@{}l}{\textbf{B. OOD：旧实现诊断，不进入主结果}}\\*[3pt]
LIBERO-Plus\par Spatial 子集 & OpenVLA-OFT & H2：七轴各 10 次；两策略同名单；保留上游提示词混杂。 & 全套件 SR 69.6\%；仅作背景参照，非本批 Spatial 子集。\hyperlink{src:PLUSNUM}{[4]} & SR 82.86\%（58/70）；诊断值，不作论文优劣比较。 \\[6pt]
LIBERO-Plus\par Spatial 子集 & AVA-VLA & H2：七轴各 10 次；两策略同名单；保留上游提示词混杂。 & NR：AVA 论文未报告本批 Plus 协议。\hyperlink{src:AVA}{[2]} & SR 78.57\%（55/70）；诊断值，不作论文优劣比较。 \\[6pt]
\multicolumn{5}{@{}l}{\textbf{C. 失败检测：固定配置，未见任务}}\\*[3pt]
SAFE rollout\par OpenVLA / WidowX & SAFE-LSTM & H3：白盒特征；3 种子；1,000 epoch；按任务划分；仅评分/校准，不干预。 & NR：原论文未列出 WidowX 域数值。\hyperlink{src:SAFE}{[5]} & 未见任务 AUROC 62.78 \(\pm\) 0.99；AUPRC 73.98 \(\pm\) 2.30。 \\[6pt]
SAFE rollout\par OpenVLA / WidowX & SAFE-MLP & H3：白盒特征；3 种子；1,000 epoch；按任务划分；仅评分/校准，不干预。 & NR：原论文未列出 WidowX 域数值。\hyperlink{src:SAFE}{[5]} & 未见任务 AUROC 81.85 \(\pm\) 3.67；AUPRC 87.34 \(\pm\) 3.84。 \\[6pt]
SAFE rollout\par OpenVLA / WidowX & Cosine kNN & H4：同 H3 划分；固定 k=5；训练轨迹建库，无梯度训练。 & NR：原论文未列出 WidowX 域数值。\hyperlink{src:SAFE}{[5]} & 未见任务 AUROC 63.64 \(\pm\) 7.51；AUPRC 67.61 \(\pm\) 7.32。 \\[6pt]
SAFE rollout\par OpenVLA / WidowX & Euclidean kNN & H4：同 H3 划分；固定 k=5；训练轨迹建库，无梯度训练。 & NR：原论文未列出 WidowX 域数值。\hyperlink{src:SAFE}{[5]} & 未见任务 AUROC 53.67 \(\pm\) 6.92；AUPRC 64.45 \(\pm\) 7.30。 \\[6pt]
SAFE rollout\par OpenVLA / WidowX & RND & H5：同划分；3 种子\(\times\)200 epoch；小尾批合并；仅检测，无恢复。 & NR：原论文未列出 WidowX 域数值。\hyperlink{src:SAFE}{[5]} & 未见任务 AUROC 49.41 \(\pm\) 4.01；AUPRC 59.56 \(\pm\) 4.96。 \\[6pt]
SAFE rollout\par Pi0-FAST / DROID & SAFE-LSTM & H3：白盒特征；3 种子；1,000 epoch；按任务划分；仅评分/校准，不干预。 & 未见任务 AUROC 58.70 \(\pm\) 4.37；Real Franka 列，论文调优配置。\hyperlink{src:SAFE}{[5]} & 未见任务 AUROC 53.58 \(\pm\) 2.08；AUPRC 53.93 \(\pm\) 4.75。 \\[6pt]
SAFE rollout\par Pi0-FAST / DROID & SAFE-MLP & H3：白盒特征；3 种子；1,000 epoch；按任务划分；仅评分/校准，不干预。 & 未见任务 AUROC 64.16 \(\pm\) 5.88；Real Franka 列，论文调优配置。\hyperlink{src:SAFE}{[5]} & 未见任务 AUROC 53.53 \(\pm\) 4.45；AUPRC 61.34 \(\pm\) 3.01。 \\[6pt]
SAFE rollout\par Pi0-FAST / DROID & Cosine kNN & H4：同 H3 划分；固定 k=5；训练轨迹建库，无梯度训练。 & 未见任务 AUROC 59.51 \(\pm\) 5.76；Real Franka 列，论文调优配置。\hyperlink{src:SAFE}{[5]} & 未见任务 AUROC 49.43 \(\pm\) 3.57；AUPRC 49.19 \(\pm\) 2.54。 \\[6pt]
SAFE rollout\par Pi0-FAST / DROID & Euclidean kNN & H4：同 H3 划分；固定 k=5；训练轨迹建库，无梯度训练。 & 未见任务 AUROC 60.27 \(\pm\) 4.79；Real Franka 列，论文调优配置。\hyperlink{src:SAFE}{[5]} & 未见任务 AUROC 53.74 \(\pm\) 3.97；AUPRC 52.60 \(\pm\) 1.51。 \\[6pt]
SAFE rollout\par Pi0-FAST / DROID & RND & H5：同划分；3 种子\(\times\)200 epoch；小尾批合并；仅检测，无恢复。 & 未见任务 AUROC 45.83 \(\pm\) 5.10；Real Franka 列，论文调优配置。\hyperlink{src:SAFE}{[5]} & 未见任务 AUROC 48.87 \(\pm\) 1.96；AUPRC 49.38 \(\pm\) 0.98。 \\[6pt]
\multicolumn{5}{@{}l}{\textbf{D. 图像/时序拼图判断：作者派生测试集}}\\*[3pt]
SMF-CALVIN\par 公开 1p 数据 & I-FailSense\par 完整模型 & H6：单视角时序拼图＋文本；203 样本；固定发布权重；离线投票/解码。 & 1p accuracy 90.64\%；Table I。\hyperlink{src:IFS}{[6]} & Accuracy 90.64\%（184/203）；与论文保留位一致。 \\[6pt]
SMF-CALVIN\par 公开 1p 数据 & I-FailSense\par LoRA/VLM 组件 & H6：单视角时序拼图＋文本；203 样本；固定发布权重；离线投票/解码。 & 1p accuracy 85.71\%；Table I。\hyperlink{src:IFS}{[6]} & Accuracy 85.71\%（174/203）；与论文保留位一致。 \\[6pt]
\multicolumn{5}{@{}l}{\textbf{E. 世界模型：公开小型验证协议}}\\*[3pt]
WCM\par pick-place val & WCM & H7：固定公开权重；18 episode / 1,617 窗口；值与下一状态预测。 & NR：未定位同一 quick-val 的误差报告；论文 RL 成功率不可替代。\hyperlink{src:WCM}{[7]} & Value RMSE 0.380342；MAE 0.196550；Pearson 0.720664；next-state MSE 0.007358。 \\[6pt]
\multicolumn{5}{@{}l}{\textbf{F. 数据生成与环境检查：不是模型性能}}\\*[3pt]
ManiSkill3\par 三任务 & 官方演示\par 动作回放 & H8：PickCube、StackCube、PegInsertionSide 各 3 条；执行动作，不逐帧写真值。 & N/A：工程回放，不是学习策略论文指标。 & 9/9 回放成功；只说明演示可执行。 \\[6pt]
ManiSkill3\par 两任务 & FailGen\par 失败注入器 & H9：grasp / trans\_x；seed 7/8；按环境终局判定，不把注入直接当失败。 & N/A：不是 AHA 检测准确率。 & 8/8 链路完成；4 成功、4 失败。 \\[6pt]
\end{longtable}\endgroup
\clearpage\section{结果解释与可改进空间}
\subsection{固定策略与提升空间}
OFT 50/50、AVA 48/50 来自同一组小样本 Spatial 初态，不能据此宣称超过论文，也不能用这一接近满分的协议作为新方法主要卖点。正式对照冻结同一执行策略；比较不加 harness、固定查询、轻量风险触发与预算自适应查询，保持底座、输入权限和恢复动作库一致。

\subsection{LIBERO-Plus 的旧分数不回填成修正后成绩}
OFT 58/70、AVA 55/70 是有效的执行诊断，但上游把扰动后缀混入部分非语言任务指令。两列分数不能与全量论文分数直接相减。正式协议对非语言轴固定标准任务文本，按预注册名单重新运行；旧证据保持原样。

\subsection{SAFE：区分调优缺口与研究缺口}
DROID 上论文 Real Franka 列的 MLP AUROC 为 64.16，本次固定配置为 53.53；这不是证明原方法失效，也不能把补齐基线调优当作新方法贡献。WidowX 未在原论文表中给出对应值。已有真实域结果支持继续研究跨任务判断，但 DROID 标签异常、特征选择与校准稳定性仍需在开发/校准划分中处理。\hyperlink{src:SAFE}{[5]}

\subsection{I-FailSense：对应成绩吻合，结论范围有限}
单视角完整模型 accuracy 90.64\%，组件 85.71\%，均与论文 Table I 的保留位一致。已有固定权重和作者派生划分使其成为最接近论文对应条目的复现。该结果不是未触碰 OOD 保证，也不含可交互换视角或恢复；论文的 success-F1 与我们记录的 failure-F1 不直接比较。\hyperlink{src:IFS}{[6]}

\subsection{WCM 与数据生成不混排}
WCM 公开验证集只含成功轨迹，值误差不能说明失败识别能力。9/9 演示回放和 FailGen 的 8 次注入用于工程验收，不能证明 ManiSkill3 已饱和、AHA 检测有效或主动恢复成功。

\subsection{与本研究的区别}
已有实验覆盖固定执行策略、被动风险检测、离线视觉判断与短时序预测。新增研究对象是“何时缺证据、查询什么、何时停止或恢复”，并由受约束程序执行决策。代码迭代仅使用训练/开发失败案例；测试反馈不进入代码修改循环。当前报告中的空白属于尚待验证的能力，不是已经获得的增益。

\clearpage\section{2026 顶会的数据集数量核查}
纳入范围：与本研究直接相关的记忆、主动感知、鲁棒性、效率及策略改进方法，覆盖 ICLR、ICML、CVPR 2026 主会。只收录有正式会议条目或官方接收信息且能核对实验设置的 8 篇。SAFE 属于 NeurIPS 2025；WCM 属于 2026 预印本参考，二者不进入这个计数样本。

统一规则：只计用于机器人任务结果评测的命名仿真 benchmark；训练语料不计。LIBERO 四/五套件合并，SimplerEnv 的 Bridge/Fractal 合并，RoboTwin 1.0/2.0 合并；有独立任务/划分的派生 benchmark 单列，但不等同于独立仿真平台。真机实验单列，不能叫作第三个公开数据集。

\begingroup
\small
\setlength{\tabcolsep}{3pt}
\begin{longtable}{@{}L{30mm}L{22mm}L{51mm}L{10mm}L{49mm}@{}}
\toprule
\textbf{论文} & \textbf{会议} & \textbf{仿真 benchmark} & \textbf{计数} & \textbf{真机与计数说明} \\ \midrule
\endfirsthead
\toprule
\textbf{论文} & \textbf{会议} & \textbf{仿真 benchmark} & \textbf{计数} & \textbf{真机与计数说明} \\ \midrule
\endhead
\bottomrule\endfoot
HAMLET \hyperlink{src:HAM}{[10]}\hyperlink{src:HAMVENUE}{[11]} & ICLR 2026 & LIBERO；RoboCasa Kitchen；SimplerEnv & 3 & 有；跨基准迁移，不只摘要中的两套 \\ \addlinespace[5pt]
MemoryVLA \hyperlink{src:MEM}{[8]}\hyperlink{src:MEMVENUE}{[9]} & ICLR 2026 & LIBERO；SimplerEnv；MIKASA-Robo & 3 & 有；Bridge/Fractal 合并，不把真机两类任务拆算 \\ \addlinespace[5pt]
SP-VLA \hyperlink{src:SP}{[12]}\hyperlink{src:SPFULL}{[13]} & ICLR 2026 & LIBERO；SimplerEnv & 2 & 有；效率与任务质量同时评测 \\ \addlinespace[5pt]
SimpleVLA-RL \hyperlink{src:SIMPLE}{[14]} & ICLR 2026 & LIBERO；RoboTwin 1.0 / 2.0 & 2 & 有；若版本分开，原文计 3 \\ \addlinespace[5pt]
Cosmos Policy \hyperlink{src:COSMOS}{[15]} & ICLR 2026 & LIBERO；RoboCasa & 2 & 有；RoboCasa 不是 RoboCasa365 \\ \addlinespace[5pt]
AVA-VLA \hyperlink{src:AVA}{[2]} & CVPR 2026 & LIBERO；CALVIN & 2 & 有；四个 LIBERO 套件只计 1 \\ \addlinespace[5pt]
ActiveVLA \hyperlink{src:ACTIVE}{[16]}\hyperlink{src:ACTIVEFULL}{[17]} & CVPR 2026 & RLBench；COLOSSEUM；GemBench & 3 & 有；三者共享 RLBench 系仿真底座 \\ \addlinespace[5pt]
StableVLA \hyperlink{src:STABLE}{[18]}\hyperlink{src:STABLEVENUE}{[19]} & ICML 2026 & LIBERO；CALVIN & 2 & 有；多个扰动等级不增加数据集数 \\ \addlinespace[5pt]
\end{longtable}
\endgroup
统一计数均值 2.375、中位数 2，8/8 落在 2-3 之间。若将 RoboTwin 两版本分开，变为 4 篇用 2 个、4 篇用 3 个，仍不能推出“大多数恰好 3 个”。MemoryVLA 原文把部分协议及真机类别细分为 6 个 benchmark，也不等于 6 个独立数据集。\hyperlink{src:MEM}{[8]}\hyperlink{src:SIMPLE}{[14]}

结论不是凑齐三套数据，而是覆盖不同失效机制，并有跨平台或真实域证据。没有实体机器人时，真实离线数据能够补充域外检测证据，但不能替代真机主动观察与恢复实验。本文研究结论将限定在已实际验证的范围。

\clearpage\section{数据集的剩余空间与取舍}
“未饱和”在这里指：在明确版本、任务和训练条件下，强基线仍留下可研究的失败空间；不是对某数据集所有任务、所有模型的永久判断。以下论文数字不是最新全社区排行榜，更不是对不同 benchmark 的横向难度排序。

\begingroup
\small
\setlength{\tabcolsep}{3pt}
\begin{longtable}{@{}L{28mm}L{54mm}L{40mm}L{44mm}@{}}
\toprule
\textbf{数据集 / 协议} & \textbf{可核验的性能证据} & \textbf{本研究定位} & \textbf{现有状态与边界} \\ \midrule
\endfirsthead
\toprule
\textbf{数据集 / 协议} & \textbf{可核验的性能证据} & \textbf{本研究定位} & \textbf{现有状态与边界} \\ \midrule
\endhead
\bottomrule\endfoot
LIBERO 标准套件 & OFT 四套件均值 97.1\%；AVA 联合策略 98.0\%；Cosmos Policy 98.5\%。\hyperlink{src:OFT}{[1]}\hyperlink{src:AVA}{[2]}\hyperlink{src:COSMOS}{[15]} & 接近上限：保留控制项，不以 clean SR 单独证明创新。 & 已有 Spatial 两策略结果；其他套件不是本次全量复现。 \\ \addlinespace[5pt]
LIBERO-Plus & 作者 OFT 全套件 69.6\%；相机 56.4\%、初态 31.9\%；作者 mix-SFT 整体 79.5\%。\hyperlink{src:PLUSNUM}{[4]} & 选入：视觉/初态扰动、证据不足与恢复主闭环。 & 代码、资产与采集接口已有；正式无提示词混杂结果待运行。 \\ \addlinespace[5pt]
MIKASA-Robo\par MemoryVLA 五任务协议 & 同表 OFT 28.4\%、MemoryVLA 41.2\%；记忆相关任务仍有明显空间。\hyperlink{src:MEM}{[8]} & 选入：独立 ManiSkill 系平台；评测历史证据检索与部分可观测验证。 & 已有 ManiSkill 基础，不等于 MIKASA 已安装/复现；需锁定旧协议、数据与权重。 \\ \addlinespace[5pt]
SAFE Pi0-FAST / DROID & 论文最佳所列 MLP 未见任务 AUROC 64.16；本次固定配置 53.53。\hyperlink{src:SAFE}{[5]} & 选入：真实域外部离线失败检测；不是 DROID 全集。 & 已有 1,464 条原始 rollout、780 条既用筛选数据；标签审计已有。 \\ \addlinespace[5pt]
RoboCasa365 & GR00T N1.5 预训练＋目标全数据微调：atomic 68.5\%、composite-seen 40.6\%、composite-unseen 42.1\%。\hyperlink{src:ROBOCASA}{[22]} & 保留扩展：明确未接近满分，但控制适配与长时序成本更高。 & 未作为已完成基线；不把旧 RoboCasa 的 67.1\% 当作 365 版成绩。 \\ \addlinespace[5pt]
SMF-CALVIN / SAFE-WidowX & I-FailSense accuracy 90.64\%；WidowX 本次 MLP AUROC 81.85。\hyperlink{src:IFS}{[6]} & 保留诊断：语义判断、跨策略检测；不额外重复凑主数据集数量。 & 已完成相应离线基线；不能据此证明交互式工具路由。 \\ \addlinespace[5pt]
\end{longtable}
\endgroup
\clearpage\section{确定的评测组合与资源边界}
主评测组合确定为“两套仿真基准＋一套真实离线数据”：LIBERO-Plus、MIKASA-Robo、SAFE-DROID；标准 LIBERO 是干净控制项。它覆盖视觉扰动、时序记忆和真实跨任务三种问题，而不是把三个指标不同的分数平均成一个总分。

MIKASA 接入采用 MemoryVLA 论文的五任务协议：ShellGameTouch、InterceptMedium、RememberColor3/5/9。论文设置为每任务 250 演示、100 评测 episode、128\(\times\)128 图像；当前官方仓库已存在 90 任务 VLA 版本，不能将 41.2\% 移植成新版本成绩。环境提交、任务 ID、控制器与数据版本一致后才进入数值主表。\hyperlink{src:MEM}{[8]}\hyperlink{src:MIKASA}{[20]}\hyperlink{src:MIKASAV}{[21]}

MIKASA 的环境和数据公开，能够复用 ManiSkill 的工程经验，但发布策略/控制适配及状态分叉仍需验收。此项是新的研发接入工作，不是零成本增加第三套数据。仅下载固定五任务所需资源，不默认获取整个新版本数据；底层 privileged state 只供标签和验收，不开放给模型。

历史证据与新观察必须分开：记忆任务中未曾记录的过去画面不能事后“购买”，当前多拍一帧也不一定恢复被遮挡的历史事实。harness 只能检索真实已记录历史，且所有对照遵循相同的缓存、保留和检索成本规则。

SAFE-DROID 只证明真实数据上的检测迁移。既有轨迹不能生成任意视角的物理反事实或真实恢复成功标签；已用 780 条也不能再称作全新未触碰测试。标签异常清理、任务划分与数据使用历史全部显式记录。

RoboCasa365 作为长时序家庭场景扩展，不在当前主线同时增加移动底盘、全套 365 任务或完整大模型重训。其微调后约 40\% 的复合任务成绩比接近零分的零样本任务更适合区分验证/恢复能力与底层策略根本不会执行的情况。\hyperlink{src:ROBOCASA}{[22]}

\clearpage\section{统一对照与论文使用条件}
\subsection{同底座、同权限}
主闭环固定执行策略。OFT、AVA、MemoryVLA 的换底座比较与同底座加 harness 比较分开。RGB、历史、深度、分割、内部特征分别标注权限；禁止把获得额外特权输入的收益写成路由本身的提升。

\subsection{固定可比较的 harness}
主对照包含：无附加验证、固定频率/固定多视角、被动 critic 触发、预算自适应证据路由。恢复库、最大动作步数与预算上限一致；增加程序迭代模块时，以同一冻结初始版本作对照，报告迭代调用与执行成本。现有 H1-H9 不等于这些对照已完成。

\subsection{按任务根隔离}
先固定训练、开发、校准、测试清单及哈希。相同初态根的扰动、观察分支与恢复分支不能跨集合。程序生成只使用开发反馈；测试时冻结代码。自演化训练预算与部署期查询预算分开记账。

\subsection{报告真正相关的结果}
闭环按数据集分别报告任务 SR、失败恢复 SR、误干预率、证据调用数、GPU 秒及 p95 延迟。检测使用相同时间窗、相同正类的 AUROC/AUPRC、固定误报水平下的检出率与提前量。拒判同时报告覆盖率，不把“少报警”当作“更安全”。

\subsection{难度与公平性验收}
在开发集检查失败是否可观察、可干预，排除基础设施错误与完全不可执行状态。难度和任务名单在最终测试前固定，不根据新方法优势筛选；统计区间按任务/episode 聚类，避免把同一轨迹的多帧当独立样本。

\subsection{本报告的完成边界}
已有结果足以作为起始基线与方法开发依据，但尚无跨上述主组合、同预算、同输入的完整新方法对照。数据集尚未饱和不意味着一定存在可由主动观察解决的剩余错误，更不意味着已获得顶会级贡献。

\clearpage\section{证据索引与参考来源}
表中所有 our reproduced performance 来自仓库已有结果，不来自论文或新模拟数字。生成脚本从原始汇总 JSON 重算 SAFE 均值/样本标准差，并保存每个输入文件 SHA-256 与表行对应关系。

H1/H2 的详细协议见 OPENVLA\_AVA\_VLA.md、LIBERO\_PLUS.md；H3-H5 见 SAFE\_20260905.md、SAFE\_DROID\_20260905.md、EXTENDED\_BASELINES\_20260905.md、SAFE\_RND\_20260905.md；H6-H9 见 IFAILSENSE\_20260905.md、WCM.md、MANISKILL3.md、AHA\_FAILGEN.md。这些历史明细保留，不因总报告替换而删除。

ActiveVLA、AHA、ADV/VeGAS、VAP-TAMP、Thea、Self-Harness/Zetta 不列为已完成数值基线。旧状态审计仍保留其日期；本次特别复查 ActiveVLA 官方入口仍为待发布，因此不能将 AVA-VLA 的结果代替 ActiveVLA。\hyperlink{src:ACTIVECODE}{[23]}

\begingroup\setlength{\parskip}{2pt}
\par\hypertarget{src:OFT}{}[1] \href{https://arxiv.org/html/2502.19645}{OpenVLA-OFT：原论文 Table I（2025）}
\par\hypertarget{src:AVA}{}[2] \href{https://openaccess.thecvf.com/content/CVPR2026/papers/Xiao_AVA-VLA_Improving_Vision-Language-Action_models_with_Active_Visual_Attention_CVPR_2026_paper.pdf}{AVA-VLA：CVPR 2026，Table 1、Table 2 与真机实验}
\par\hypertarget{src:PLUS}{}[3] \href{https://openaccess.thecvf.com/content/CVPR2026/html/Fei_LIBERO-Plus_A_Progressive_Robustness_Benchmark_for_Visual-Language-Action_Models_CVPR_2026_paper.html}{LIBERO-Plus：CVPR 2026；原论文结果表}
\par\hypertarget{src:PLUSNUM}{}[4] \href{https://arxiv.org/pdf/2510.13626}{LIBERO-Plus：作者论文完整表与各套件明细}
\par\hypertarget{src:SAFE}{}[5] \href{https://arxiv.org/html/2506.09937v2}{SAFE：NeurIPS 2025，Table 8、B.8 与真实 Franka 实验}
\par\hypertarget{src:IFS}{}[6] \href{https://arxiv.org/html/2509.16072v1}{I-FailSense：Table I 与评测定义}
\par\hypertarget{src:WCM}{}[7] \href{https://github.com/sylvestf/WCM}{WCM：作者公开 quick benchmark、数据与权重}
\par\hypertarget{src:MEM}{}[8] \href{https://arxiv.org/html/2508.19236v2}{MemoryVLA：ICLR 2026；4.1-4.5，Table 4}
\par\hypertarget{src:MEMVENUE}{}[9] \href{https://shihao1895.github.io/MemoryVLA/}{MemoryVLA：作者 ICLR 2026 项目与公开资源}
\par\hypertarget{src:HAM}{}[10] \href{https://arxiv.org/html/2510.00695}{HAMLET：ICLR 2026；Tables 2-3、真机实验}
\par\hypertarget{src:HAMVENUE}{}[11] \href{https://iclr.cc/virtual/2026/poster/10010110}{HAMLET：ICLR 2026 正式会议条目}
\par\hypertarget{src:SP}{}[12] \href{https://proceedings.iclr.cc/paper_files/paper/2026/hash/4072543747a14bbed76284cf2c04b9e9-Abstract-Conference.html}{SP-VLA：ICLR 2026 正式会议论文}
\par\hypertarget{src:SPFULL}{}[13] \href{https://arxiv.org/html/2506.12723}{SP-VLA：实验设置与真机补充}
\par\hypertarget{src:SIMPLE}{}[14] \href{https://proceedings.iclr.cc/paper_files/paper/2026/file/cbfbcb4da14235bd69b134070898ae9d-Paper-Conference.pdf}{SimpleVLA-RL：ICLR 2026，3.1 实验设置}
\par\hypertarget{src:COSMOS}{}[15] \href{https://proceedings.iclr.cc/paper_files/paper/2026/hash/748becc400a57c0e31cfe6a2e7951467-Abstract-Conference.html}{Cosmos Policy：ICLR 2026 正式会议条目}
\par\hypertarget{src:ACTIVE}{}[16] \href{https://openaccess.thecvf.com/content/CVPR2026/html/Liu_ActiveVLA_Injecting_Active_Perception_into_Vision-Language-Action_Models_for_Precise_3D_CVPR_2026_paper.html}{ActiveVLA：CVPR 2026 正式会议论文}
\par\hypertarget{src:ACTIVEFULL}{}[17] \href{https://arxiv.org/html/2601.08325v1}{ActiveVLA：4.1 与三套仿真基准}
\par\hypertarget{src:STABLE}{}[18] \href{https://dagroup-pku.github.io/StableVLA/}{StableVLA：ICML 2026；Table 1 与真机设置}
\par\hypertarget{src:STABLEVENUE}{}[19] \href{https://icml.cc/Downloads/2026}{ICML 2026 官方论文列表（含 StableVLA）}
\par\hypertarget{src:MIKASA}{}[20] \href{https://github.com/CognitiveAISystems/MIKASA-Robo}{MIKASA-Robo：官方数据、环境与版本说明}
\par\hypertarget{src:MIKASAV}{}[21] \href{https://github.com/CognitiveAISystems/MIKASA-Robo/releases}{MIKASA-Robo：90 任务 VLA 版与旧 RL 版发布边界}
\par\hypertarget{src:ROBOCASA}{}[22] \href{https://arxiv.org/html/2603.04356}{RoboCasa365：Table 2，目标任务微调后的成功率}
\par\hypertarget{src:ACTIVECODE}{}[23] \href{https://github.com/ZhenyangLiu/ActiveVLA-Injecting-Active-Perception-into-VLA}{ActiveVLA：本次核查仍待发布的代码与评测入口}
\endgroup
\end{document}